% Part: proof-theory
% Chapter: proof-search
% Section: completeness

\documentclass[../../../include/open-logic-section]{subfiles}

\begin{document}

\olfileid{pt}{ps}{cpl}

\olsection{সম্পূর্ণতা}

এবার দেখাই, প্রমাণ-অন্বেষণ
অ্যালগরিদমটি নির্ভরযোগ্য:
সেটি ব্যর্থ হয়ে থেমে গেলে
বা কখনো না থামলে শুরুর
সিকোয়েন্ট~$\Gamma \Sequent \Delta$-র
কোনো !!{proof} নেই। এ জন্য
!!a{structure}~$\Struct M$ গড়ি,
যাতে প্রতিটি $!A \in \Gamma$-র জন্য
$\Sat{M}{!A}$ এবং প্রতিটি
$!A \in \Delta$-র জন্য
$\Sat/{M}{!A}$।
অর্থাৎ~$\Struct M$-এ
$\Gamma$-র সব !!{formula}s সত্য,
$\Delta$-র সব !!{formula}s মিথ্যা;
তাই $\Gamma \Sequent \Delta$
বৈধ নয়। \Log{G3c} যথার্থ
বলে $\Gamma \Sequent \Delta$-র
কোনো !!{proof} নেই।
% BN-SRC-591 TARGET-CORRECTION-BEGIN
\par\smallskip\noindent\textnormal{\small\textbf{উৎস-সংশোধন (BN-SRC-591):} মিথ্যা উত্তরাংশের শর্তে উৎসের $! \in \Delta$ অসম্পূর্ণ সূত্র; পরের $\Sat/{M}{!A}$-এর সঙ্গে মিলিয়ে $!A \in \Delta$ করা হয়েছে।}
% BN-SRC-591 TARGET-CORRECTION-END

!!{formula}s-এর জোড়া
$\Theta, \Xi$ এবং
!!{constant}s-এর সেট~$C$ থেকে
$\Struct{M}$ সংজ্ঞায়িত করি।
এই $\Theta, \Xi$ ও~$C$
প্রমাণ-অন্বেষণের
ব্যর্থ (থেমে যাওয়া বা অনন্তকাল চলা)
প্রক্রিয়ার একটি
\emph{ব্যর্থতা-শাখা} থেকে পাই।
এমন শাখায় সিকোয়েন্টের অনুক্রম
$\Pi_n \Sequent \Lambda_n$ ও
!!{constant}s-এর সেট~$C_n$ থাকে।
এখানে $\Pi_0 \Sequent \Lambda_0$
হলো শুরুর সিকোয়েন্ট
$\Gamma \Sequent \Delta$;
আর শাখার পরবর্তী ধাপ থাকলে
$\Pi_{n+1} \Sequent \Lambda_{n+1}$
হলো $n+1$-তম পর্যায়ে তৈরি
$\Pi_n \Sequent \Lambda_n$-এর
এমন পূর্বধারণা, যার !!a{proof}
অ্যালগরিদম পায় না। শাখার
প্রতিটি চলমান~$n$-এ
এমন পূর্বধারণা থাকতেই হবে;
নইলে অ্যালগরিদম !!a{proof} পেত।
$C_n$ হলো
$\Pi_n \Sequent \Lambda_n$-র
সঙ্গে যুক্ত !!{constant}s-এর সেট।
% BN-SRC-592 TARGET-CORRECTION-BEGIN
\par\smallskip\noindent\textnormal{\small\textbf{উৎস-সংশোধন (BN-SRC-592):} শেষ বাক্যের সিকোয়েন্টে উৎসের $\Delta_n$ চিহ্নটি এই শাখায় সংজ্ঞায়িত নয়; শাখার ডান অংশ $\Lambda_n$ নেওয়া হয়েছে।}
% BN-SRC-592 TARGET-CORRECTION-END
% BN-SRC-596 TARGET-CORRECTION-BEGIN
\par\smallskip\noindent\textnormal{\small\textbf{উৎস-সংশোধন (BN-SRC-596):} উৎসে প্রতিটি $n$-এ পরবর্তী ব্যর্থ পূর্বধারণা দাবি করার পরেই থেমে-যাওয়া সসীম শাখাকেও ব্যর্থতা-শাখা বলা হয়েছে। তাই পরবর্তী-ধাপের শর্ত চলমান পর্যায়গুলিতে সীমিত করা হয়েছে।}
% BN-SRC-596 TARGET-CORRECTION-END

অ্যালগরিদম যদি কেবল
পারমাণবিক !!{formula}s-যুক্ত
একটি সর্বোচ্চ সিকোয়েন্টে
থেমে যায়, সেখানে শেষ হওয়া
শাখাটি ব্যর্থতা-শাখা।
আর অ্যালগরিদম কখনো না থামলে
ক্রমে বড় বড় বৃক্ষ তৈরি করে।
এগুলিকে একটি অনন্ত বৃক্ষের
বৃদ্ধিশীল পর্যায় ভাবতে পারি
(অসীম সংখ্যক ধাপ শেষে
যে বৃক্ষ পাওয়া যেত)।
বৃক্ষটি অনন্ত হলেও
সসীম-শাখাবিশিষ্ট: প্রতিটি
নোডের অব্যবহিত উপরে
একটি বা দুটি সন্তান-সিকোয়েন্ট।
K\H{o}nig-এর লেমা অনুযায়ী
প্রতিটি অনন্ত সসীম-শাখাবিশিষ্ট
বৃক্ষের অনন্ত শাখা আছে।
অন্বেষণ অনন্তকাল চললে
এমন অনন্ত শাখাই ব্যর্থতা-শাখা।

$\Pi_n \Sequent \Lambda_n$
ও $C_n$-এর অনুক্রম
ব্যর্থতা-শাখা হলে ধরি
$\Theta = \bigcup_n \Pi_n$,
$\Xi = \bigcup_n \Lambda_n$
(!!{formula}s-এর সূচক
ও একাধিক সংঘটন সরিয়ে)
এবং $C = \bigcup_n C_n$।

এবার~$\Struct{M}$ সংজ্ঞায়িত করতে পারি।
\begin{enumerate}
\item $\Domain{M}$ হলো~$C$-র
ধ্রুবক ও
$\Gamma \Sequent \Delta$-তে থাকা
!!{function}s দিয়ে গড়া,
কিন্তু কোনো !!{variable}s না-থাকা
পদগুলির সেট।
\item $c \in C$ হলে
$\Assign{c}{M} = c$।
% BN-SRC-595 TARGET-CORRECTION-BEGIN
\par\smallskip\noindent\textnormal{\small\textbf{উৎস-সংশোধন (BN-SRC-595):} ধ্রুবকের ব্যাখ্যায় $c$-কে উৎসে সমগ্র পদ-ডোমেনের সদস্য ধরা হয়েছে; $\Assign{c}{M}$ কেবল ধ্রুবক-চিহ্নের জন্য, তাই $c \in C$ করা হয়েছে।}
% BN-SRC-595 TARGET-CORRECTION-END
\item $f$ যদি
$\Gamma \Sequent \Delta$-তে
$m$-স্থানীয় !!{function} হয়
এবং $t_1$, \dots,
$t_m \in \Domain{M}$,
তবে $\Assign{f}{M}(t_1, \dots, t_m) =
f(t_1, \dots, t_m)$।
% BN-SRC-593 TARGET-CORRECTION-BEGIN
\par\smallskip\noindent\textnormal{\small\textbf{উৎস-সংশোধন (BN-SRC-593):} $m$-স্থানীয় ফাংশনের ফলপদে উৎসের শেষ যুক্তি $t_n$; আগের যুক্তিতালিকা অনুযায়ী $t_m$ করা হয়েছে।}
% BN-SRC-593 TARGET-CORRECTION-END
\item $R$ যদি
$\Gamma \Sequent \Delta$-তে
$m$-স্থানীয় !!{predicate} হয়,
তবে $\Assign{R}{M} =
\Setabs{\tuple{t_1, \dots, t_m} \in
\Domain{M}^m}{R(t_1, \dots, t_m) \in \Theta}$।
% BN-SRC-594 TARGET-CORRECTION-BEGIN
\par\smallskip\noindent\textnormal{\small\textbf{উৎস-সংশোধন (BN-SRC-594):} $m$-স্থানীয় বিধেয়ের ব্যাখ্যায় উৎসের ডোমেন-ঘাত ও শেষ যুক্তি $n$ হলেও বাঁ টিউপল $m$-দৈর্ঘ্যের; সংজ্ঞা ও পরের ব্যাখ্যায় সর্বত্র $m$ নেওয়া হয়েছে।}
% BN-SRC-594 TARGET-CORRECTION-END
\end{enumerate}
$\Struct{M}$ একটি
\emph{পদ-মডেল}: তার ডোমেন
পদের সেট, আর !!{constant}s
ও !!{function}s-এর ব্যাখ্যা
এমন যে $\Value{t}{M} = t$।
$\Theta$-র পারমাণবিক
!!{formula}s দিয়ে
$\Assign{R}{M}$ এমনভাবে
সংজ্ঞায়িত যে
$\Sat{M}{R(t_1, \dots, t_m)}$
তখনই এবং কেবল তখনই,
যখন $R(t_1, \dots, t_m)
\in \Theta$।

\begin{lem}\ollabel{lem:truth}
প্রতিটি $!A \in \Theta \cup \Xi$-র জন্য:
\begin{enumerate}
  \item $!A \in \Theta$ হলে $\Sat{M}{!A}$।
  \item $!A \in \Xi$ হলে $\Sat/{M}{!A}$।
\end{enumerate}
\end{lem}

\begin{proof}
  $\depth{!A}$-র উপর আরোহ করি।

প্রথমে ধরি $!A$ পারমাণবিক,
অর্থাৎ $!A \ident \lfalse$ অথবা
$!A \ident R(t_1, \dots, t_m)$।

যেহেতু $\Pi_n \Sequent \Lambda_n$
ব্যর্থতা-শাখায় আছে,
কোনো $\Pi_n$-তে
$\lfalse$ থাকতে পারে না
(থাকলে $\Pi_n \Sequent \Lambda_n$
স্বতঃসিদ্ধ হতো)।
$\Struct{M}$-এর সংজ্ঞা থেকে
$\Sat{M}{R(t_1, \dots, t_m)}$
তখনই এবং কেবল তখনই,
যখন $R(t_1, \dots, t_m) \in
\Theta$। ফলে $!A \in
\Theta$ হলে $\Sat{M}{!A}$।

$!A$ পারমাণবিক হলে,
$!A \in \Pi_n$ থেকে
$!A \in \Pi_{n+1}$,
আর $!A \in \Lambda_n$ থেকে
$!A \in \Lambda_{n+1}$ হয়।
(অন্বেষণ-অ্যালগরিদমের
উত্তরসূরি সিকোয়েন্ট তৈরির
পদ্ধতি থেকে এটি মেলে।)
তাই $!A \in \Theta \cap \Xi$
হলে কোনো~$n$-এর জন্য
$!A \in \Pi_n \cap
\Lambda_n$ হতো। তখন
$\Pi_n \Sequent \Lambda_n$
স্বতঃসিদ্ধ হতো; কিন্তু
ব্যর্থতা-শাখায় স্বতঃসিদ্ধ
থাকতে পারে না।
অতএব পারমাণবিক~$!A$-এর
জন্য $!A \notin \Theta
\cap \Xi$। ফলে
$!A \in \Xi$ হলে
$!A \notin \Theta$;
$\Struct{M}$-এর সংজ্ঞা থেকে
$!A \in \Xi$ হলে
$\Sat/{M}{!A}$।

আরোহী ধাপে ধরি,
$!A$-র চেয়ে কম
!!{depth}-এর সব
!!{formula}s-এর জন্য
(1) ও (2) সত্য।
ক্ষেত্রভেদ করে $!A$-র
জন্যও (1) ও (2) দেখাই।

প্রথমে লক্ষ করি,
$\Pi_n \Sequent \Lambda_n$-তে
$!A^i$ থাকলে, $i$
সেখানকার ক্ষুদ্রতম সূচক না হওয়া
পর্যন্ত $!A^i$ পরের
$\Pi_{n+1} \Sequent \Lambda_{n+1}$-এও
থাকে। প্রতি পর্যায়ে যোগ হওয়া
সর্বোচ্চ সিকোয়েন্টে ক্ষুদ্রতম
সূচক বাড়ে; তাই অবশেষে
এমন $\Pi_n \Sequent \Lambda_n$
পাই যেখানে $!A^i$ আছে
এবং $i$ ক্ষুদ্রতম।
$n+1$-তম পর্যায়ে অ্যালগরিদম
$!A^i$ হ্রাস করে।
অর্থাৎ, $!A \in \Theta$
হলে কোনো $n$-এ
$\Pi_n = \Pi_n', !A^i$
এবং $i$ ক্ষুদ্রতম সূচক।
একইভাবে, $!A \in \Xi$
হলে কোনো $n$-এ
$\Lambda_n = \Lambda_n', !A^i$
এবং $i$ ক্ষুদ্রতম সূচক।

\begin{enumerate}
  \item $!A \in \Theta$
  এবং $!A \ident !B \land !C$।
  তখন কোনো $n$-এর জন্য
  $\Pi_n = \Pi_n', (!B \land !C)^i$।
  $\Pi_{n+1} \Sequent
  \Lambda_{n+1}$ হলো
  \LeftR{\land}-এর সংশ্লিষ্ট
  পূর্বধারণা, অর্থাৎ
  $\Pi_{n+1} = \Pi_n', !B^k, !C^{k+1}$।
  ফলে $!B, !C \in
  \Theta$। আরোহ-অনুমান থেকে
  $\Sat{M}{!B}$ ও $\Sat{M}{!C}$।
  $\Sat{M}{}$-এর সংজ্ঞা থেকে
  $\Sat{M}{!B \land !C}$।

  \item $!A \in \Xi$
  এবং $!A \ident !B \land !C$।
  তখন কোনো $n$-এর জন্য
  $\Lambda_n = \Lambda_n', (!B \land !C)^i$।
  $\Pi_{n+1} \Sequent
  \Lambda_{n+1}$ হলো
  \RightR{\land}-এর
  দুটি পূর্বধারণার একটি;
  সিদ্ধান্ত
  $\Pi_n \Sequent \Lambda'_n, (!B \land !C)^i$।
  অর্থাৎ, $\Lambda_{n+1} =
  \Lambda_n', !B^k$
  অথবা $\Lambda_{n+1} =
  \Lambda_n', !C^k$।
  ফলে $!B \in \Xi$
  অথবা $!C \in \Xi$।
  আরোহ-অনুমান থেকে
  $\Sat/{M}{!B}$ অথবা
  $\Sat/{M}{!C}$;
  $\Sat{M}{}$-এর সংজ্ঞা থেকে
  $\Sat/{M}{!B \land !C}$।
  % BN-SRC-597 TARGET-CORRECTION-BEGIN
  \par\smallskip\noindent\textnormal{\small\textbf{উৎস-সংশোধন (BN-SRC-597):} ডান-∧ ক্ষেত্রের দুই সিকোয়েন্টে হ্রাসযোগ্য সূত্রের $i$ সূচক উৎসে বাদ ছিল; অ্যালগরিদম ও বাম-∧ ক্ষেত্র অনুযায়ী $(!B\land!C)^i$ ফিরেছে।}
  % BN-SRC-597 TARGET-CORRECTION-END

  \item $!A \ident \lnot !B$,
  $!A \ident !B \lor !C$
  ও $!A \ident !B \lif !C$
  ক্ষেত্রগুলি একইরকম;
  অনুশীলনে রাখা হলো।

  \item $!A \in \Theta$
  এবং $!A \ident \lexists[x][!B(x)]$।
  তখন কোনো $n$-এ
  $\Pi_n = \Pi_n', \lexists[x][!B(x)]^i$।
  $\Pi_{n+1}
  \Sequent \Lambda_{n+1}$
  হলো \LeftR{\lexists}-এর
  সংশ্লিষ্ট পূর্বধারণা,
  অর্থাৎ কোনো~$c$-এর জন্য
  $\Pi_{n+1} = \Pi_n', !B(c)^k$।
  ফলে $!B(c) \in \Theta$
  এবং~$c \in C$।
  আরোহ-অনুমান থেকে
  $\Sat{M}{!B(c)}$;
  $\Sat{M}{}$-এর সংজ্ঞা থেকে
  $\Sat{M}{\lexists[x][!B(x)]}$।

  \item $!A \in \Xi$
  এবং $!A \ident \lexists[x][!B(x)]$।
  তখন কোনো $n$-এ
  $\Lambda_n = \Lambda_n', \lexists[x][!B(x)]^i$।
  প্রতিবার $\lexists[x][!B(x)]$
  হ্রাসের পরে পূর্বধারণার
  ডান দিকে $\lexists[x][!B(x)]$
  নতুন সূচকসহ
  থেকে যায়। তাই ব্যর্থতা-শাখার
  ডানে থাকলে
  $\lexists[x][!B(x)]$
  অসীমবার হ্রাস পায়।
  প্রতিটি $t \in
  \Domain{M}$ একসময়
  ``$C_n$-র ধ্রুবকগুলি
  দিয়ে গড়া এবং
  ডানে $\lexists[x][!B(x)]$
  হ্রাসে এখনও ব্যবহৃত নয়
  এমন প্রথম পদ'' হয়।
  তাই প্রতিটি $t \in \Domain{M}$-র
  জন্য কোনো~$n$-এ
  $!B(t) \in \Lambda_n$,
  এবং ফলে $!B(t) \in \Xi$।
  আরোহ-অনুমান থেকে
  সব $t \in \Domain{M}$-র
  জন্য $\Sat/{M}{!B(t)}$;
  $\Sat{M}{}$-এর সংজ্ঞা থেকে
  $\Sat/{M}{\lexists[x][!B(x)]}$।

  \item $!A \ident \lforall[x][!B(x)]$
  ক্ষেত্রগুলি একইরকম;
  অনুশীলনে রাখা হলো।
\end{enumerate}
\end{proof}

\begin{cor}\ollabel{cor:complete}
\Log{G3c}-র প্রমাণ-অন্বেষণ
অ্যালগরিদম সম্পূর্ণ:
সেটি ব্যর্থ হয়ে থেমে গেলে
বা কখনো না থামলে শুরুর
সিকোয়েন্ট~$\Gamma \Sequent \Delta$
বৈধ নয়, তাই তার কোনো
!!{proof} নেই।
\end{cor}

\begin{proof}
  অ্যালগরিদম ব্যর্থ হয়ে
  থেমে গেলে বা কখনো
  না থামলে একটি
  ব্যর্থতা-শাখা আছে,
  যা থেকে $\Struct{M}$
  সংজ্ঞায়িত করা যায়।
  \olref{lem:truth} থেকে
  সব $!A \in \Gamma$-র
  জন্য $\Sat{M}{!A}$
  এবং সব $!A \in \Delta$-র
  জন্য $\Sat/{M}{!A}$।
  (মনে রাখি,
  $\Pi_0 = \Gamma$ ও
  $\Lambda_0 =
  \Delta$; তাই
  $\Gamma \subseteq \Theta$
  এবং $\Delta \subseteq \Xi$।)
  অতএব $\Gamma \Sequent \Delta$
  বৈধ নয়। \Log{G3c}-র
  যথার্থতা থেকে
  $\Gamma \Sequent \Delta$-র
  কোনো !!{proof} নেই।
\end{proof}

\begin{cor}
\Log{G3c} সম্পূর্ণ:
$\Gamma \Sequent \Delta$
বৈধ হলে
$\Log{G3c} \Proves
\Gamma \Sequent \Delta$।
\end{cor}

\begin{proof}
  বিপরীতার্থক রূপ প্রমাণ করি।
  ধরি, $\Log{G3c} \Proves/
  \Gamma \Sequent \Delta$।
  খুঁজে পাওয়ার মতো
  কোনো !!{proof} নেই;
  তাই প্রমাণ-অন্বেষণ
  অ্যালগরিদম ব্যর্থ হয়ে
  থামবে অথবা কখনো থামবে না।
  \olref{cor:complete} থেকে
  $\Gamma \Sequent \Delta$
  বৈধ নয়।
\end{proof}


\end{document}
